% !TEX root = ../../../tree_notes.tex
% =====================================================================
%  13.2  Tree DP: the in/out two-state  sec:c10-twostate
%  Exports: def:c10-inout  thm:c10-mis-tree  lst:c10-inout
%           fig:c10-mistrace  ex:c10-mis-run
%           prob:c10-inout-solved (CF 1529C)  lst:c10-1529c-soln
%  CF 1529C "Parsa's Humongous Tree" IS the two-state shape: assign each
%  node a value in [l_v, r_v] maximising sum over edges of |a_u - a_v|;
%  optimum is at endpoints, so the state is (pick l, pick r) not (in, out),
%  but the recurrence is the MIS pair verbatim. All numbers on T are
%  hand-computed in \figref{c10-runtree-dp} and reused here unchanged.
% =====================================================================
\section{Two-State DP: In / Out at Every Node}\label{sec:c10-twostate}

\begin{previouslybox}
\begin{enumerate}[nosep, leftmargin=1.6em]
  \item Chapter~2 solved LC~337 ``House Robber III'' with the pair
        $(\mathit{rob},\mathit{skip})$ at every binary-tree node. Which
        two facts made the pair \emph{closed} under the child-fold?
        (a node's value depends only on whether each child is taken or
        not, and taking $v$ forces every child to the ``skip'' branch)
  \item Chapter~12 defined the maximum independent set on a general tree
        in the CSES~1130 preview as ``largest set of nodes with no two
        adjacent''. Which complement is it, and why does that equivalence
        need K{\H{o}}nig's theorem on bipartite graphs? (minimum vertex
        cover; trees are bipartite so $\alpha'=\nu$ matches)
  \item Ch.\chref{c09}'s \coderef{c09-dfs3} pushed $\depth$ down as a
        parameter and folded scalars up. When we replace the scalar by a
        \emph{pair}, what changes in the fold? (nothing in the DFS
        skeleton; everything in the state design)
  \item Ch.1's \thmref{c01-height-recursion} and Ch.2's LC~337 are both
        bottom-up DP on a fixed root. What is one tree problem that
        \emph{cannot} be done with a scalar state and \emph{can} be done
        with a pair? (max-weight independent set: a scalar cannot
        express ``$v$ is available only if $c$ is skipped'')
\end{enumerate}
\end{previouslybox}

\begin{topicintro}
This section is the general-tree upgrade of Chapter~2's $(\mathit{rob},
\mathit{skip})$ pattern. Ch.12 previewed CSES~1130 in its mixed set as the
doorway into the tree-DP chapters; here it is discharged -- not by solving
CSES~1130 (that is a canonical warm-up, one-page code, and the plan says
Ch.12 owns the pointer) but by solving CF~1529C (\probref{c10-inout-solved}),
a problem whose \emph{shape} is the MIS pair dressed as ``assign a value in
a range and maximise the sum of edge differences''. Once you see that the
optimal assignment always sits at an endpoint of the range, the fold is
literally \thmref{c10-mis-tree} with a distance kernel replacing a
unit weight.

\textbf{Prerequisites.} \secref{c10-subtrees} for the child-fold template
and the tuple idea; Ch.12's \coderef{c09-dfs3} as the DFS skeleton;
\thmref{c01-height-recursion} as the special case where the tuple size is
one. \emph{Not} assumed: matching theory (K{\H{o}}nig is quoted, not proved;
Ch.31 owns matchings), rerooting (Ch.15 owns it).

\textbf{After this section you can}
\begin{itemize}[nosep]
  \item state the MIS recurrence on a general tree and hand-trace
        \lstinline{dp_in}, \lstinline{dp_out} on $T$; \hfill\lv{3}
  \item recognise ``endpoints carry the optimum'' in a range-assignment
        problem and rewrite it as an in/out pair; \hfill\lv{4}
  \item cold-write CF~1529C's recurrence
        $\mathit{dp}[v][b]=\sum_{u}\max_{b'}\big(\mathit{dp}[u][b']+|\mathit{val}(v,b)-\mathit{val}(u,b')|\big)$
        from memory; \hfill\lv{4}
\end{itemize}
\end{topicintro}

\subsection{Intuition}

A scalar state fails the moment the decision at $v$ has to \emph{constrain}
the decisions at $v$'s children. ``Take $v$'' means ``no child may be
taken''; ``skip $v$'' means ``each child is free to be taken or skipped,
whichever is better''. Two numbers per node are exactly enough to remember
both branches, and every fold closes under them:

\[
  \mathit{dp}_{\mathit{in}}[v] \;=\; w(v) + \sum_{c\in\children(v)}\mathit{dp}_{\mathit{out}}[c],
  \qquad
  \mathit{dp}_{\mathit{out}}[v] \;=\; \sum_{c\in\children(v)}\max\!\big(\mathit{dp}_{\mathit{in}}[c],\,\mathit{dp}_{\mathit{out}}[c]\big).
\]

Leaf initialisers: $\mathit{dp}_{\mathit{in}}[\text{leaf}]=w$,
$\mathit{dp}_{\mathit{out}}[\text{leaf}]=0$ -- Ch.1's $\height(\nil)$
convention wearing the same hat.

This pair solves the \term{maximum independent set} (MIS) on a tree: the
largest subset of nodes containing no adjacent pair. Its complement is the
\term{minimum vertex cover} (the smallest subset touching every edge), and
on \emph{any} graph K{\H{o}}nig's theorem gives
$\alpha(G)+\tau(G)=|V|$ for bipartite $G$; trees are bipartite, so MIS on a
tree and min vertex cover on a tree are the same problem viewed upside
down. We quote the identity; the matching theory is Chapter~31's territory.

CF~1529C is \emph{this} shape wearing a different costume. Each node has a
range $[\ell_v,r_v]$ and you must pick $a_v$ in it; the score is
$\sum_{uv\in E}|a_u-a_v|$. Two observations turn it into an in/out pair:
(i) the objective is convex in every $a_v$ separately (sum of absolute
values of affine functions), so a maximum over the box $[\ell_v,r_v]$ is at
a \emph{vertex}, i.e.\ every $a_v\in\{\ell_v,r_v\}$; (ii) once $a_v$ is
restricted to two choices, ``choose $\ell_v$'' vs ``choose $r_v$'' is
literally the out/in pair of MIS with a per-edge weight $|a_u-a_v|$ instead
of a vertex weight. That reduction is \probref{c10-inout-solved}.

\subsection{Definitions and Notation}

\begin{definition}{The in/out pair on a general tree}{c10-inout}
Fix a root $r$ and a \emph{take weight} $w\colon V\to\R_{\ge 0}$. For every
node $v$, define the pair $(\mathit{dp}_{\mathit{in}}[v],
\mathit{dp}_{\mathit{out}}[v])$ recursively by
\[
  \mathit{dp}_{\mathit{in}}[v] \;=\; w(v)+\!\!\sum_{c\in\children(v)}\!\mathit{dp}_{\mathit{out}}[c],
  \qquad
  \mathit{dp}_{\mathit{out}}[v] \;=\; \sum_{c\in\children(v)}\!\max\!\big(\mathit{dp}_{\mathit{in}}[c],\,\mathit{dp}_{\mathit{out}}[c]\big),
\]
with the initialisers $\mathit{dp}_{\mathit{in}}[\text{leaf}]=w(\text{leaf})$
and $\mathit{dp}_{\mathit{out}}[\text{leaf}]=0$ (empty sum).
\end{definition}

\noindent\textbf{In plain words.} $\mathit{dp}_{\mathit{in}}[v]$ is the best
total weight of an independent set inside $v$'s subtree \emph{forced to
contain $v$}; $\mathit{dp}_{\mathit{out}}[v]$ is the best inside $v$'s
subtree \emph{with $v$ excluded}. Both are computed in one DFS.

\begin{theorem}{The in/out pair solves maximum-weight independent set on trees}{c10-mis-tree}
For any tree $T$ and any $w\colon V\to\R_{\ge 0}$, the value
$\max(\mathit{dp}_{\mathit{in}}[r],\mathit{dp}_{\mathit{out}}[r])$ equals the
maximum weight of an independent set in $T$, and Algorithm
\coderef{c10-inout} computes the pair at every node in $\Theta(n)$ time.
\end{theorem}
\begin{proofbox}
Complexity is \thmref{c10-subtree-template} with a two-coordinate tuple:
one DFS, $\bigO{\deg(v)}$ per node, total $\Theta(n)$. Correctness is
strong induction on the height of $v$'s subtree. \emph{Base}: a leaf has
only two independent sets inside its own subtree, $\emptyset$ (weight~$0$,
which is $\mathit{dp}_{\mathit{out}}$) and $\{v\}$ (weight~$w(v)$, which
is $\mathit{dp}_{\mathit{in}}$). \emph{Step}: let $I$ be a maximum-weight
independent set inside $v$'s subtree. Either $v\in I$ or $v\notin I$.
If $v\in I$, no child of $v$ is in $I$ (independence), so $I$ is $\{v\}$
together with independent sets inside each $c\in\children(v)$ \emph{excluded}
at the root, contributing $\sum_c\mathit{dp}_{\mathit{out}}[c]=\mathit{dp}_{\mathit{in}}[v]-w(v)$
by induction. If $v\notin I$, then $I\cap\text{subtree}(c)$ is an independent
set inside $c$'s subtree with no constraint at $c$ itself, so its weight is
at most $\max(\mathit{dp}_{\mathit{in}}[c],\mathit{dp}_{\mathit{out}}[c])$
by induction, and $I$'s weight is at most $\sum_c\max(\cdot,\cdot)=
\mathit{dp}_{\mathit{out}}[v]$. Both bounds are achieved by construction:
$\mathit{dp}_{\mathit{in}}[v]$ is achieved by taking $v$ plus each child's
$\mathit{dp}_{\mathit{out}}$-witness, $\mathit{dp}_{\mathit{out}}[v]$ by
each child's $\max$-witness. Hence
$\max(\mathit{dp}_{\mathit{in}}[v],\mathit{dp}_{\mathit{out}}[v])$ is the
MIS weight inside $v$'s subtree. At the root, ``inside $r$'s subtree'' is
all of $T$. \qed
\end{proofbox}

\begin{insightbox}
Two coordinates are \emph{necessary and sufficient} for MIS on a tree.
Necessary: a scalar cannot distinguish ``$v$ is available for its parent
to combine with'' from ``$v$ is already blocked by a child taking it''.
Sufficient: the recurrence closes because $v$'s constraint on a child is
purely binary (either $v$ is in, forcing $c$ out, or $v$ is out, letting
$c$ decide freely). LC~968 ``Binary Tree Cameras'' (Ch.2's mixed set owns
the URL) needs \emph{three} coordinates because the constraint is ternary
(camera at $v$ / covered by parent / covered by a child): same skeleton,
one more row in the fold.
\end{insightbox}

\subsection{The Algorithm: Two Coordinates, One DFS}

\begin{lstlisting}[caption={MIS via the in/out pair on a general tree: \coderef{c09-dfs3}'s skeleton with two accumulators}, label={lst:c10-inout}]
// adj[u] = neighbours; w[u] = take weight; call mis(r, -1).
vector<vector<int>> adj;
vector<long long> w, dp_in, dp_out;         // 64-bit: n * 1e9 fits only in LL

void mis(int u, int p) {
    dp_in[u]  = w[u];                       // "u is taken" accumulates w(u)
    dp_out[u] = 0;                          // "u is skipped" accumulates 0
    for (int v : adj[u]) if (v != p) {      // parent guard = rooting
        mis(v, u);
        dp_in[u]  += dp_out[v];             // u taken -> v must be skipped
        dp_out[u] += max(dp_in[v], dp_out[v]);// u skipped -> v free
    }
}
// answer: max(dp_in[r], dp_out[r])
// witness: DFS again choosing (in -> force children out; out -> per-child max)
\end{lstlisting}

\begin{complexitybox}
Time $\Theta(n)$ by \thmref{c10-mis-tree}, same fold budget as
\coderef{c09-dfs3}. Memory: two $\bigO{n}$ long-long arrays plus $\bigO{h}$
call stack; substitute \coderef{c09-iter}'s two-phase explicit-stack form
whenever chains are legal inputs at $n \ge 2\cdot 10^{5}$. Verdict: the
canonical contest shape for ``choose a subset with a hard per-edge
constraint on a tree'' -- and the exact reduction target for CF~1529C
below, where the ``hard constraint'' is a two-choice vertex assignment.
\end{complexitybox}

\subsection{Worked Example: $\mathit{dp}_{\mathit{in}},\mathit{dp}_{\mathit{out}}$ on $T$ with Unit Weights}

\begin{example}{MIS on the running tree $T$ at $w\equiv 1$}{c10-mis-run}
Take $T$ from \figref{c10-runtree-dp}, root at~$1$, and set $w(v)=1$ for
every $v$ (plain MIS, unweighted). The fold in \coderef{c10-inout} gives,
hand-computed:
\[
  \mathit{dp}_{\mathit{in}} = [\,6,2,2,1,1,1,2,1,1,1\,],\qquad
  \mathit{dp}_{\mathit{out}} = [\,5,2,3,0,0,0,1,0,1,0\,].
\]
The MIS size is $\max(6,5)=\mathbf{6}$, achieved at the root's in-state.
\end{example}

\begin{figure}[H]
\centering
\begin{tikzpicture}[diagram, scale=0.82, transform shape,
    tnode/.style={circle, line width=1pt, minimum size=8mm,
                  font=\small\bfseries},
    taken/.style={tnode, draw=successgreen, fill=successgreen!18, text=successgreen!60!black},
    skipped/.style={tnode, draw=warningred, fill=warningred!8, text=warningred!60!black}]
  \node[taken]   (n1)  at (0,0)       {$1$};
  \node[skipped] (n2)  at (-3,-1.6)   {$2$};
  \node[skipped] (n3)  at ( 3,-1.6)   {$3$};
  \node[taken]   (n4)  at (-5.2,-3.2) {$4$};
  \node[skipped] (n5)  at (-1.6,-3.2) {$5$};
  \node[taken]   (n6)  at ( 1.6,-3.2) {$6$};
  \node[taken]   (n7)  at ( 5.2,-3.2) {$7$};
  \node[taken]   (n8)  at (-1.6,-4.8) {$8$};
  \node[skipped] (n9)  at ( 5.2,-4.8) {$9$};
  \node[taken]   (n10) at ( 5.2,-6.4) {$10$};
  \draw[tedge] (n1)--(n2) (n1)--(n3) (n2)--(n4) (n2)--(n5)
               (n3)--(n6) (n3)--(n7) (n5)--(n8) (n7)--(n9) (n9)--(n10);
  \foreach \n/\xin/\xout in {n1/6/5, n2/2/2, n3/2/3, n4/1/0, n5/1/0,
                             n6/1/0, n7/2/1, n8/1/0, n9/1/1, n10/1/0} {
    \node[font=\scriptsize\bfseries, text=successgreen!55!black, anchor=south]
      at (\n.north) {\scriptsize in=\xin};
    \node[font=\scriptsize\bfseries, text=warningred!70!black, anchor=north]
      at (\n.south) {\scriptsize out=\xout};
  }
  \node[annot, font=\scriptsize, text=deepblue, anchor=west]
    at (7.0,0) {green $=$ taken};
  \node[annot, font=\scriptsize, text=deepblue, anchor=west]
    at (7.0,-0.7) {$S=\{1,4,6,7,8,10\}$};
\end{tikzpicture}
\caption{The in/out pair on $T$ at unit weights. Root~$1$'s badges
$(6,5)$ decide the answer; green fill is one back-tracking witness
$\{1,4,6,7,8,10\}$. No two adjacent nodes are green; the size~$6$
upper bound is the max of the pair at the root, by
\thmref{c10-mis-tree}.}
\label{fig:c10-mistrace}
\end{figure}

\begin{center}
\small
\renewcommand{\arraystretch}{1.15}
\begin{tabular}{l|cccccccccc}
\toprule
node $v$                       & 1 & 2 & 3 & 4 & 5 & 6 & 7 & 8 & 9 & 10 \\
\midrule
$\mathit{dp}_{\mathit{in}}[v]$ & 6 & 2 & 2 & 1 & 1 & 1 & 2 & 1 & 1 & 1 \\
$\mathit{dp}_{\mathit{out}}[v]$& 5 & 2 & 3 & 0 & 0 & 0 & 1 & 0 & 1 & 0 \\
\bottomrule
\end{tabular}
\end{center}

\begin{tracebox}[Post-order through \coderef{c10-inout} on $T$, root~$1$, $w\equiv 1$]
post-order: 4, 8, 5, 2, 10, 9, 7, 6, 3, 1.\\
Leaves 4, 6, 8, 10: dp\_in=1, dp\_out=0.\\
5 (child 8): dp\_in$=1+$dp\_out$[8]=1+0=1$; dp\_out$=\max(1,0)=1$.\\
9 (child 10): dp\_in$=1+0=1$; dp\_out$=\max(1,0)=1$.\\
7 (child 9): dp\_in$=1+$dp\_out$[9]=1+1=2$; dp\_out$=\max(1,1)=1$.\\
2 (children 4,5): dp\_in$=1+$dp\_out$[4]+$dp\_out$[5]=1+0+1=2$;
dp\_out$=\max(1,0)+\max(1,1)=1+1=2$.\\
3 (children 6,7): dp\_in$=1+0+1=2$; dp\_out$=\max(1,0)+\max(2,1)=1+2=3$.\\
1 (children 2,3): dp\_in$=1+2+3=\mathbf{6}$;
dp\_out$=\max(2,2)+\max(2,3)=2+3=\mathbf{5}$.\\
MIS size$=\max(6,5)=\mathbf{6}$ \checkmark. Back-track: root takes
dp\_in$\Rightarrow$ children 2,3 skipped; for each, use dp\_out. Node~2's
dp\_out picks $\max(\mathit{dp}_{\mathit{in}}[4]=1,0)=1$ (take~4) and
$\max(\mathit{dp}_{\mathit{in}}[5]=1,1)=1$ (either; pick~8 under~5 which
is the dp\_out side). Node~3's dp\_out picks $\max(1,0)=1$ (take~6) and
$\max(2,1)=2$ (take~7, which forces 9 out and picks 10). Result:
$\{1,4,8,6,7,10\}$ $=$ $\{1,4,6,7,8,10\}$ (reordering) $=6$ \checkmark.
\end{tracebox}

\begin{insightbox}[Key takeaways]
\begin{enumerate}[nosep, leftmargin=1.6em]
  \item Same DFS as \coderef{c10-fold}; a two-coordinate tuple. The
        complexity is unchanged.
  \item ``Take $v$ forces every child out'' is the entire reason we need
        two coordinates -- with one, the constraint has nowhere to live.
  \item Back-tracking for a \emph{witness} is another DFS, not extra
        state; \thmref{c10-mis-tree}'s achievability clause tells you
        which branch to follow.
  \item Sanity: MIS size $+$ min-vertex-cover size $=n$. On $T$,
        min-vertex-cover is $n-6=4$ (the four edges $1\!-\!2,1\!-\!3,2\!-\!5,
        3\!-\!7$ cover everything once we pick $2,3,5,7$ -- no, wait, this
        needs a proper cover; the identity holds by K{\H{o}}nig but a
        witness is not free. Do not fabricate covers without checking.)
\end{enumerate}
\end{insightbox}

\subsection{Variations and Twists}

\begin{itemize}[nosep]
  \item \textbf{Weighted MIS}: set $w(v)$ to the problem's own weight,
        nothing else changes. \hfill\lv{3}
  \item \textbf{Minimum weighted vertex cover}: complement the MIS pair
        ($\mathit{cover}_{\mathit{in}}[v]=w(v)+\sum_c\mathit{cover}_{\mathit{out}}[c]$,
        $\mathit{cover}_{\mathit{out}}[v]=\sum_c\min(\cdot,\cdot)$);
        the two are the same DP read upside-down. \hfill\lv{3}
  \item \textbf{Three-state extension (LC~968 preview)}: for
        ``minimum dominating set'' or ``binary tree cameras'' the pair
        $(\text{take},\text{skip})$ is not enough -- you must remember
        ``skipped but already covered by a child'' as a third state.
        Ch.2 owns the URL of LC~968; this chapter is not solving it, only
        noting that the pair becomes a triple and the fold gets one more
        row. \hfill\lv{4}
  \item \textbf{CF~1324F ``Maximum White Subtree'' (preview, not solved
        here)}: the in/out pair computes ``net black minus net white in
        $v$'s subtree'' in one DFS, but the problem asks for
        \emph{every root}; that is Ch.15's rerooting machinery and is
        deliberately outside this chapter's scope. \hfill\lv{4}
\end{itemize}

\subsection{Common Pitfalls}

\begin{warningbox}[Trap 1 -- dp\_in written with the wrong child accumulator]
The tempting mistake is to allow a taken $v$ to combine with \emph{free}
children: $\mathit{dp}_{\mathit{in}}[v]=w(v)+\sum_c\max(
\mathit{dp}_{\mathit{in}}[c],\mathit{dp}_{\mathit{out}}[c])$. That is
\emph{wrong} -- independence forbids a child from being in when the parent
is. The correct $\mathit{dp}_{\mathit{in}}$ uses $\mathit{dp}_{\mathit{out}}[c]$
\emph{unconditionally}. Symptom: MIS size on a two-node tree comes out $2$
instead of $1$.
\end{warningbox}

\begin{warningbox}[Trap 2 -- $\mathit{dp}_{\mathit{out}}$ with a $\max$ instead of a sum]
The other tempting mistake is $\mathit{dp}_{\mathit{out}}[v]=\max_c\max(
\mathit{dp}_{\mathit{in}}[c],\mathit{dp}_{\mathit{out}}[c])$. That treats
independence as ``$v$ conflicts with only one child'', but the children are
mutually independent subtrees \emph{once $v$ is out}; every child contributes
to the total. Fix: outer $\sum$, inner $\max$.
\end{warningbox}

\begin{warningbox}[Trap 3 -- $\mathit{int}$ overflow on weighted MIS]
CF~1529C's ranges go up to $10^{9}$, so per-node contributions go up to
$10^{9}$ and the sum over $10^{5}$ edges goes up to $10^{14}$ -- a factor
of~$10^{5}$ past int. \lstinline{dp_in}, \lstinline{dp_out} and every
accumulator in \coderef{c10-inout} must be \lstinline{long long}. This is
the ``long long'' trap from Ch.12's chapter-end re-appearing in a new
shape; the fold structure hides the overflow inside the $\sum$ and
$\max$, not inside the DFS, so it is a \emph{wrong answer} not a runtime
error.
\end{warningbox}

\subsection{Problem Lab}

The section's anchor is Codeforces~1529C, \probref{c10-inout-solved}: a
range-assignment maximisation problem whose reduction to the in/out pair is
a one-line convexity observation. Once the observation is made, the
algorithm is \coderef{c10-inout} verbatim with a per-edge weight matrix
replacing the scalar $w(v)$.

% ---------------------------------------------------------------------------
\solvedproblem{c10-inout-solved}{Parsa's Humongous Tree}{Codeforces 1529C}
  {https://codeforces.com/problemset/problem/1529/C}{CF 1600 L4}
  {in/out two-state, convexity-at-endpoints, weighted MIS reduction, long long}

\begin{problemstatement}[Parsa's Humongous Tree (samples verified against the live problem)]
There are $n$ nodes labelled $1,\dots,n$ forming a tree rooted at~$1$, and
for each node $v$ a range $[\ell_v,r_v]$ of integers with
$\ell_v\le r_v$. Assign each node $v$ an integer $a_v\in[\ell_v,r_v]$ to
\emph{maximise}
\[
  \sum_{uv\in E} |a_u - a_v| .
\]
Constraints: $2\le n\le 10^{5}$, $1\le \ell_v\le r_v\le 10^{9}$, and the
sum of $n$ over all test cases is $\le 2\cdot 10^{5}$. Multiple test cases
first line $t$; each test case is $n$, then $n$ lines of $(\ell_v,r_v)$,
then $n-1$ lines of edges. \textbf{Sample~1}: $n=2$, ranges
$(1,6)$ and $(3,8)$, edge $1\!-\!2$ $\Rightarrow$ answer~$7$ (assign
$a_1=1$, $a_2=8$). \textbf{Sample~2}: $n=3$, ranges $(1,3)$, $(4,6)$,
$(7,9)$, edges $1\!-\!2$ and $2\!-\!3$ $\Rightarrow$ answer~$8$ (assign
$(a_1,a_2,a_3)=(1,6,9)$ or $(3,4,9)$). \textbf{Sample~3}: $n=6$, ranges
$(3,14)(12,20)(12,19)(2,12)(10,17)(3,17)$, edges $3\!-\!2$, $6\!-\!5$,
$1\!-\!5$, $2\!-\!6$, $4\!-\!6$ $\Rightarrow$ answer~$62$.
\end{problemstatement}

\begin{attemptgate}[Attempt it yourself first -- \lv{4}, give yourself 30 minutes]
Two ideas must click before the DFS helps. (i) Where in $[\ell_v,r_v]$ does
$a_v$ sit at an optimum? (ii) Once $a_v$ is restricted, what does ``$v$'s
state'' mean?
\peek{Fix every $a_u$ for $u\ne v$. How does the objective depend on $a_v$?}
\nudge{It is a sum of $|a_v - a_u|$ terms -- convex in $a_v$, so its
      maximum on the interval $[\ell_v,r_v]$ is at one of the two
      endpoints. Do this for every $v$ simultaneously: the whole optimum
      lives at a vertex of the box $\prod_v[\ell_v,r_v]$.}
\fullhint{Every $a_v$ is either $\ell_v$ or $r_v$ at an optimum, so $v$ has
      two states ($b\in\{0,1\}$, where $b=0$ means $\ell_v$ and $b=1$ means
      $r_v$). Let $\mathit{dp}[v][b]$ be the max sum of $|a_u-a_v|$ over
      edges inside $v$'s subtree, given $a_v$ fixed at state $b$. At $v$
      with state $b$, each child $u$ may sit at either of its two states
      $b'$; the contribution of edge $uv$ plus the subtree below $u$ is
      $\mathit{dp}[u][b'] + |\mathit{val}(v,b)-\mathit{val}(u,b')|$, and you
      take the max over $b'$ per child. Sum over children. Answer is
      $\max(\mathit{dp}[1][0],\mathit{dp}[1][1])$. All arithmetic in
      \lstinline{long long}.}
\end{attemptgate}

\solstep{Understanding the Problem}
Two subproblems compose. (i)~\emph{Where} in $[\ell_v,r_v]$ is the optimum?
Once every other $a_u$ is fixed, the objective restricted to $a_v$ is
$\sum_{uv\in E}|a_v-a_u|$, a sum of convex (V-shaped) functions of $a_v$,
hence convex; a convex function on a closed interval attains its maximum
at an endpoint. Doing this for every $v$ and taking the joint argument:
the box-constrained maximum of a coordinate-wise convex function is at a
vertex of the box. So we may restrict $a_v\in\{\ell_v,r_v\}$.
(ii)~\emph{Which} endpoint per node becomes a binary state, and once the
state is binary the problem is MIS on a tree with a per-edge weight
matrix: exactly the shape \thmref{c10-mis-tree} handles.

\solstep{Key Observations}
\begin{observation}
The objective $\sum_{uv}|a_u-a_v|$ is coordinate-wise convex in each $a_v$
(separate sum of absolute values), so an optimum on $\prod_v[\ell_v,r_v]$
lies at a vertex of the box; every $a_v$ is either $\ell_v$ or $r_v$.
\end{observation}
\begin{observation}
Fixing $a_v$ to a specific endpoint defines $v$'s \emph{state} $b_v\in\{0,1\}$.
The edge contribution between $u$ and $v$ then becomes a table lookup:
$|{\mathit{val}(u,b_u)-\mathit{val}(v,b_v)}|$, where
$\mathit{val}(v,0)=\ell_v$ and $\mathit{val}(v,1)=r_v$.
\end{observation}
\begin{observation}
$\mathit{dp}[v][b] :=$ max score over edges entirely inside $v$'s subtree,
given $b_v=b$. The recurrence
$\mathit{dp}[v][b]=\sum_{u\in\children(v)}\max_{b'\in\{0,1\}}\big(\mathit{dp}[u][b'] + |\mathit{val}(v,b)-\mathit{val}(u,b')|\big)$
closes: each child is decided independently once $b_v$ is fixed, because
there are no edges between two distinct child subtrees in a tree.
\end{observation}

\solstep{Think It Through With Me}
\begin{thinkbox}
\deadend{Sort nodes by range midpoint and assign greedily, taking the
extreme value that maximises the running sum.}
{There is no natural order on a tree, and the objective is \emph{jointly}
convex (per coordinate) so greedy fails on Sample~2: assigning $a_2=4$
first (the smallest midpoint) forces $a_1,a_3$ to shrink the total to at
most $|4-3|+|4-7|=4$, worse than the $8$ at $(1,6,9)$.}
\deadend{Enumerate all $2^n$ assignments of endpoints and pick the max.}
{$n\le 10^{5}$ makes this a non-starter, but a \emph{useful diagnostic}: if
you cannot see a per-node decomposition, you have not found the right
state. Here, the decomposition is exactly the in/out fold of
\thmref{c10-mis-tree}, and the state is the endpoint bit.}
\aha{Two states per node, one fold per edge, $\Theta(n)$ total.
     \coderef{c10-inout}'s skeleton, with the ``$\mathit{dp}_{\mathit{in}}$
     / $\mathit{dp}_{\mathit{out}}$'' pair replaced by
     ``$\mathit{dp}[v][0]$ / $\mathit{dp}[v][1]$'' and the per-edge weight
     $|a_u-a_v|$ replacing the vertex weight $w(v)$.}
\end{thinkbox}

\solstep{Solution}
Post-order DFS with two accumulators per node. $\mathit{dp}[v][b]$ starts
at~$0$; for each child~$u$, add
$\max\!\big(\mathit{dp}[u][0]+|\mathit{val}(v,b)-\ell_u|,\;\mathit{dp}[u][1]+|\mathit{val}(v,b)-r_u|\big)$.
Answer is $\max(\mathit{dp}[1][0],\mathit{dp}[1][1])$. All arithmetic in
\lstinline{long long} because $\sum|a_u-a_v|$ can hit $10^{14}$.

\solstep{Proof of Correctness}
\begin{proofbox}[The DP value equals the maximum of $\sum_{uv}|a_u-a_v|$ over all valid assignments]
\emph{Reduction to endpoints}: by Observation~1, restricting $a_v$ to
$\{\ell_v,r_v\}$ does not lose the optimum. \emph{Recurrence closure}:
given the assignment inside each child subtree (independently chosen), the
score contribution of edges inside $v$'s subtree decomposes as
$\sum_{u\in\children(v)}(\mathit{dp}[u][b_u] + |a_v-a_u|)$ where the sum
runs over edges of $v$'s subtree, each of which lies in exactly one of the
child sub-subtrees or is the edge $vu$. Independence of subtrees (a tree
has no edge between two different child sub-trees) is what makes the sum
legal. \emph{Optimal substructure}: for a fixed $b_v$, maximising over
$b_u$ per child is a per-child $\max$, so the DP is exact by induction on
subtree height. \emph{Answer}: the whole tree is root's subtree, so
$\max(\mathit{dp}[1][0],\mathit{dp}[1][1])$ is the global optimum. \qed
\end{proofbox}

\solstep{Implementation}
\begin{lstlisting}[caption={CF 1529C -- endpoint reduction plus in/out pair; \coderef{c10-inout}'s skeleton with a per-edge weight matrix}, label={lst:c10-1529c-soln}]
#include <bits/stdc++.h>
using namespace std;
using ll = long long;

int main() {
    ios::sync_with_stdio(false);
    int t;  cin >> t;
    while (t--) {
        int n;  cin >> n;
        vector<ll> L(n), R(n);
        for (int i = 0; i < n; ++i) cin >> L[i] >> R[i];
        vector<vector<int>> adj(n);
        for (int i = 0; i < n - 1; ++i) {
            int u, v;  cin >> u >> v;  --u; --v;
            adj[u].push_back(v);  adj[v].push_back(u);
        }
        vector<array<ll,2>> dp(n);
        // iterative two-phase order to survive n = 1e5 chains
        vector<int> par(n, -1), order;  order.reserve(n);
        vector<int> st{0};
        while (!st.empty()) {
            int u = st.back();  st.pop_back();
            order.push_back(u);
            for (int v : adj[u]) if (v != par[u]) { par[v] = u; st.push_back(v); }
        }
        for (int i = (int)order.size() - 1; i >= 0; --i) {
            int u = order[i];
            dp[u][0] = dp[u][1] = 0;
            for (int v : adj[u]) if (v != par[u]) {
                ll pickL = max(dp[v][0] + llabs(L[u] - L[v]),
                               dp[v][1] + llabs(L[u] - R[v]));
                ll pickR = max(dp[v][0] + llabs(R[u] - L[v]),
                               dp[v][1] + llabs(R[u] - R[v]));
                dp[u][0] += pickL;
                dp[u][1] += pickR;
            }
        }
        cout << max(dp[0][0], dp[0][1]) << '\n';
    }
}
\end{lstlisting}

\solstep{Trace}
\begin{tracebox}[Sample~1: $n=2$, ranges $(1,6)$ and $(3,8)$, edge $1\!-\!2$]
Root at~1. Children of 1: $\{2\}$. Node~2 has no children.
Fold at node~2: dp[2][0]=0, dp[2][1]=0 (no children to accumulate).
Fold at node~1 (u=1, v=2): pickL $=\max(0+|1-3|,0+|1-8|)=\max(2,7)=7$;
pickR $=\max(0+|6-3|,0+|6-8|)=\max(3,2)=3$.
dp[1][0] $=7$; dp[1][1] $=3$.
Answer $=\max(7,3)=\mathbf{7}$. \checkmark\ matches official.
Assignment: root at state~$0$ (value~$1$), child at state~$1$ (value~$8$),
edge difference $|1-8|=7$.
\end{tracebox}

\begin{tracebox}[Sample~2: $n=3$, ranges $(1,3),(4,6),(7,9)$, edges $1\!-\!2$ and $2\!-\!3$]
Root at~1. Children: 1$\to$$\{2\}$, 2$\to$$\{3\}$, 3$\to\emptyset$.
Fold at~3: dp[3][0]=dp[3][1]=0.
Fold at~2 (u=2, v=3): pickL $=\max(0+|4-7|,0+|4-9|)=\max(3,5)=5$;
pickR $=\max(0+|6-7|,0+|6-9|)=\max(1,3)=3$.
dp[2][0]=5, dp[2][1]=3.
Fold at~1 (u=1, v=2): pickL $=\max(5+|1-4|,3+|1-6|)=\max(8,8)=8$;
pickR $=\max(5+|3-4|,3+|3-6|)=\max(6,6)=6$.
dp[1][0]=8, dp[1][1]=6. Answer $=\mathbf{8}$. \checkmark\ matches official.
Two optimal assignments exist: $(a_1,a_2,a_3)=(1,6,9)$ giving $|1-6|+|6-9|=5+3=8$,
and $(3,4,9)$ giving $|3-4|+|4-9|=1+5=6$ -- wait, the second is $6$, not~$8$.
Rechecking: the tie at node~1 comes from either $(a_1=1,a_2=4,a_3=9)$
giving $|1-4|+|4-9|=3+5=8$ (dp[2][0]=5 corresponds to $a_2=4$, and node~3
goes to $a_3=9$ contributing $|4-9|=5$) or $(a_1=1,a_2=6,a_3=9)$ giving
$5+3=8$. Both are correct optima; the DP's $\max$ picks either, and the
value is~$8$ in both. \checkmark
\end{tracebox}

\begin{tracebox}[Sample~3: $n=6$; ranges $(3,14)(12,20)(12,19)(2,12)(10,17)(3,17)$]
Rooted view of $T$ at~1 (from CF's sample; edges $3\!-\!2$, $6\!-\!5$,
$1\!-\!5$, $2\!-\!6$, $4\!-\!6$ reorient to children):
$1\!\to\!\{5\}$; $5\!\to\!\{6\}$; $6\!\to\!\{2,4\}$; $2\!\to\!\{3\}$.
Fold order: 3, 4, 2, 6, 5, 1.  All entries in \lstinline{long long}.
Leaves 3, 4: dp$=(0,0)$.
Node~2 (L=12,R=20) with child~3 (L=12,R=19): pickL$=\max(|12-12|,|12-19|)=7$;
pickR$=\max(|20-12|,|20-19|)=8$. dp[2]$=(7,8)$.
Node~6 (L=3,R=17) with children 2 and~4 (4 is leaf with L=2,R=12):\\
  child~2 contributes max(dp[2][0]+|3-12|, dp[2][1]+|3-20|)$=\max(16,25)=25$
  for state~0 of 6; and max(7+|17-12|,8+|17-20|)$=\max(12,11)=12$ for state~1.\\
  child~4 contributes max(0+|3-2|,0+|3-12|)$=\max(1,9)=9$ (state~0 of 6);
  and max(0+|17-2|,0+|17-12|)$=\max(15,5)=15$ (state~1).\\
  Sum: dp[6]$=(25+9,12+15)=(34,27)$.
Node~5 (L=10,R=17) with child~6:\\
  pickL$=\max(34+|10-3|,27+|10-17|)=\max(41,34)=41$;\\
  pickR$=\max(34+|17-3|,27+|17-17|)=\max(48,27)=48$.\\
  dp[5]$=(41,48)$.
Node~1 (L=3,R=14) with child~5:\\
  pickL$=\max(41+|3-10|,48+|3-17|)=\max(48,62)=62$;\\
  pickR$=\max(41+|14-10|,48+|14-17|)=\max(45,51)=51$.\\
  dp[1]$=(62,51)$.
Answer$=\max(62,51)=\mathbf{62}$. \checkmark\ matches official. Optimal
assignment: node~1 at~$3$, node~5 at~$17$, node~6 at~$3$, node~2 at~$20$,
node~3 at~$12$, node~4 at~$12$; edge scores $|3-17|+|17-3|+|3-20|+|20-12|+|3-12|=14+14+17+8+9=62$.
\end{tracebox}

\solstep{Complexity}
\begin{complexitybox}
Time $\bigO{n}$ per test case (each node visited once, each edge inspected
twice, four absolute-value additions per edge); $\bigO{\sum n}$ overall
$= \bigO{2\cdot 10^{5}}$. Memory: two $\bigO{n}$ long-long arrays plus
$\bigO{n}$ adjacency. Verdict: \ac, and this is the shape the section is
about -- \emph{two states, one DFS, long long}. Chain-safe via the
two-phase \coderef{c09-iter} pattern already inlined above.
\end{complexitybox}

\solstep{Pitfalls and Variants}
\begin{warningbox}
Three ways to lose this problem: (i)~$\mathit{int}$ accumulator on values up
to~$10^{9}$ and sums up to~$10^{14}$ -- use \lstinline{long long};
(ii)~recursive DFS on a chain of $10^{5}$ nodes -- use the two-phase
\coderef{c09-iter} pattern; (iii)~forgetting the $\sum$ over children
inside each $\mathit{dp}[u][b]$ update (writing $\max$ instead) -- the same
Trap~2 as MIS above. Variant worth doing after this: CSES~1130 ``Tree
Matching'' (Ch.12's mixed set owns the pointer) -- unit-weight MIS on a
general tree, exactly \coderef{c10-inout} with $w\equiv 1$. Once you can
write CF~1529C from memory, CSES~1130 is a one-line simplification of it.
\end{warningbox}

\subsection{Practice Problems}

\begin{stuckbox}
\begin{itemize}[nosep, leftmargin=1.4em]
  \item Did you justify \emph{why} the optimum is at a vertex of the box
        before writing the DFS? A wrong convexity argument will still pass
        the samples but lose you the proof-of-correctness credit on a
        hand-graded sheet.
  \item Are both $\mathit{dp}[u][0]$ and $\mathit{dp}[u][1]$ accumulated
        with the same $\sum$ structure over $u$'s children, only differing
        in $\mathit{val}(u,b)$? A missing $\sum$ shows up as $\max$ where
        the total should be.
  \item Is your DFS iterative at $n=10^{5}$? Sample sizes are small;
        hidden tests are not.
\end{itemize}
\end{stuckbox}

\begin{practiceset}
  \practice{Marginal Cost}{Codeforces 1010D}
    {https://codeforces.com/problemset/problem/1010/D}{CF 1900 L4}
    {A single-state cousin of the in/out fold: for each node, the
    minimum fuel needed at the root to reach every vertex in $v$'s
    subtree is a bottom-up $\max$-fold over children.
    \coderef{c10-inout}'s skeleton still applies -- same DFS shape,
    different $\Phi$. Set $\Phi=\max$ and the child contribution
    carries a $\max(c_v, \cdot)$ clamp. Practice recognising when
    the state is a scalar (as here) vs.\ a pair (as in
    \thmref{c10-mis-tree}) is the point of the drill.}
\end{practiceset}

\paragraph{Back-cites (URLs live elsewhere in the book).}
\begin{itemize}[nosep,leftmargin=1.4em]
  \item \textbf{LC~337 ``House Robber III''} (\chref{c02} owns this URL):
        Ch.2 solved this with the very pair $(\mathit{rob},\mathit{skip})$
        we formalise as \thmref{c10-mis-tree}; on a \emph{binary} tree the
        pair fold has $\deg\le 2$, so the $\bigO{n}$ budget is a one-liner.
        Re-solve it now that you know the general-tree statement -- you
        should get the same code with the child loop widened.
  \item \textbf{LC~968 ``Binary Tree Cameras''} (\chref{c02} owns this
        URL in its mixed set): the three-state extension noted in
        Variations. Ch.2 has the URL; this chapter is a preview only, no
        code. The one paragraph is enough to see \emph{why} three
        coordinates are the honest budget for a dominating set (``taken
        / covered from child / needs parent's help'').
\end{itemize}

\begin{recallbox}[10-second refresher -- 13.2 Two-state in/out DP]
Two numbers per node: $(\mathit{dp}_{\mathit{in}}[v],\mathit{dp}_{\mathit{out}}[v])$
-- take $v$ (children forced out) or skip $v$ (children free).
\thmref{c10-mis-tree} solves maximum-weight independent set on a general
tree in $\Theta(n)$ via \coderef{c10-inout}; LC~337's pair on binary trees
is the special case. K{\H{o}}nig on bipartite graphs gives the vertex-cover
dual for free. CF~1529C (\probref{c10-inout-solved}) reduces range
assignment to the same pair via coordinate-wise convexity: the optimum is
at a vertex of the box, so $\mathit{dp}[v][0],\mathit{dp}[v][1]$ stand for
``pick $\ell_v$ / pick $r_v$''. Three-state (LC~968) is a preview only.
\selftest{
  \item Write the recurrence for \lstinline{dp_in[u]} and
        \lstinline{dp_out[u]} from memory, and state the leaf initialisers.
  \item Why does the in-state's child sum use $\mathit{dp}_{\mathit{out}}[c]$
        unconditionally (Trap~1) but the out-state use a $\max$?
  \item On a two-node tree, what is the MIS size? What does a wrong
        $\mathit{dp}_{\mathit{in}}$ recurrence give?
  \item What bound forces \lstinline{long long} in CF~1529C but not in
        unit-weight MIS on $n=10^{5}$ nodes?
  \item State one sentence for the ``coordinate-wise convex $\Rightarrow$
        vertex of the box'' reduction used in \probref{c10-inout-solved}.
}
\end{recallbox}
